% ===========================================================================
% The Entropy-Compression Law Scales: Predicting Optimal Layer-wise
% Compression of 7B+ Language Models via Von Neumann Entanglement Entropy
%
% Paper 3 in the Entropy-Compression Law series
%
% Nicolas Calderon Gonzalez
% Independent Research, 3Pi Agency R&D, Bogota, Colombia
%
% Target: arXiv preprint -- cs.LG (primary), quant-ph (cross-list)
% ===========================================================================
\documentclass[11pt,a4paper]{article}

% -- Packages ---------------------------------------------------------------
\usepackage[utf8]{inputenc}
\usepackage[T1]{fontenc}
\usepackage{lmodern}
\usepackage[margin=2.5cm]{geometry}
\usepackage{amsmath,amssymb,amsfonts}
\usepackage{graphicx}
\usepackage{booktabs}
\usepackage[dvipsnames]{xcolor}
\usepackage{hyperref}
\usepackage{cleveref}
\usepackage{microtype}
\usepackage{float}
\usepackage{enumitem}
\usepackage{multirow}
\usepackage[numbers,sort&compress]{natbib}

\hypersetup{
  colorlinks=true,
  linkcolor=NavyBlue,
  citecolor=OliveGreen,
  urlcolor=BrickRed
}

% -- Custom commands --------------------------------------------------------
\newcommand{\Dmin}{D_{\min}}
\newcommand{\Dmax}{D_{\max}}
\newcommand{\Svn}{S_1}
\newcommand{\Srenyi}{S_2}
\newcommand{\nats}{\;\mathrm{nats}}
\newcommand{\ie}{\textit{i.e.}}
\newcommand{\eg}{\textit{e.g.}}

% -- Title ------------------------------------------------------------------
\title{%
  \textbf{The Entropy-Compression Law Scales:\\[4pt]
  Predicting Optimal Layer-wise Compression of 7B+ Language Models\\[2pt]
  via Von Neumann Entanglement Entropy}}

\author{
  Nicolas Calderon Gonzalez \\[2pt]
  \textit{Independent Research, 3Pi Agency R\&D} \\
  \textit{Bogot\'a, Colombia} \\
  \texttt{nicolasc@trespimedios.co}
}

\date{July 2026}

% ===========================================================================
\begin{document}
\maketitle

% ===========================================================================
% ABSTRACT
% ===========================================================================
\begin{abstract}
Every layer of a large language model has a different tolerance for
compression, but current methods discover this tolerance empirically,
one layer at a time.  We show that a single number already encodes
it: the von~Neumann entanglement entropy $\Svn$ of the tensorized
weight matrix.  In two prior papers, we established the
Entropy-Compression Law on models up to 2.7B parameters.
Here we validate the law at 7B scale: on Mistral-7B (224 matrices),
$\Svn$ predicts the minimum bond dimension for a given Frobenius
error with $R^2 = 0.997$ at $\varepsilon = 50\%$; on Qwen2-7B (196
matrices), $R^2 = 0.993$.  The law \emph{improves} with model size
because 7B models exhibit wider entropy variation across layers,
giving the regression more leverage.  We then demonstrate that the
law is practically useful: entropy-guided rank allocation
outperforms uniform allocation at every compression budget tested on
GPT-2 Small (5--38\% perplexity improvement), and the advantage
grows from 5\% to 11\% after 200 steps of LoRA healing.  As a
predictor of layer-wise compressibility, $\Svn$ dominates the
competing AlphaPruning metric ($R^2 = 0.997$ vs.\ $0.022$; the two
quantities are essentially uncorrelated, $r = -0.16$).  All results
are produced by \texttt{entropy-lens}, an open-source tool that
computes $\Svn$ for any HuggingFace model in ${\sim}45$ minutes on
an Apple~M1 laptop with no GPU.
\end{abstract}

\vspace{0.5em}
\noindent\textbf{Keywords:}
entanglement entropy, LLM compression, SVD, rank allocation,
tensor networks, matrix product states, Mistral, Qwen

% ===========================================================================
% 1. INTRODUCTION
% ===========================================================================
\section{Introduction}
\label{sec:intro}

\subsection{The rank allocation problem}
\label{sec:intro-problem}

Compressing a large language model by low-rank factorization
means replacing each weight matrix $W \in \mathbb{R}^{m \times n}$
with a product $U \Sigma V^\top$ truncated at some rank $r < \min(m,n)$.
The smaller $r$ is, the fewer parameters we store and the faster
inference runs.  But different weight matrices tolerate different
amounts of truncation.  An attention key projection that learns sparse
compatibility patterns may survive at rank 200; a feed-forward gate
projection that stores dense factual associations may need rank 2000.

The practical consequence is that \emph{uniform} rank allocation ---
assigning the same $r$ to every layer --- wastes parameters on
layers that could have been compressed further and starves layers
that needed more capacity.  The field knows this: recent adaptive
methods~\citep{Sun2025ara, Wang2024svd} search for per-layer ranks
by evaluating perplexity at many candidate configurations.  The
search is expensive, model-specific, and provides no structural
insight into \emph{why} one layer needs more rank than another.

We propose a different approach.  Instead of searching for the right
rank, we \emph{predict} it from a single structural measurement of
the weight matrix: its \emph{von~Neumann entanglement entropy}
$\Svn$.

\subsection{The Entropy-Compression Law}
\label{sec:intro-law}

In two companion papers~\citep{calderon2026cartography,
calderon2026theory}, we established the \emph{Entropy-Compression
Law}: reshape a weight matrix into a high-order tensor, decompose it
as a Matrix Product State (MPS) via sequential SVD, and measure
$\Svn$ at the dominant bipartition.  Then the minimum bond dimension
$\Dmin$ needed to stay within relative Frobenius error $\varepsilon$
obeys
\begin{equation}
  \label{eq:law}
  \Dmin(\varepsilon) \approx c(\varepsilon)\, e^{\,\alpha(\varepsilon)\,\Svn}\,,
\end{equation}
where $c(\varepsilon)$ and $\alpha(\varepsilon)$ are fit constants
that depend on the error tolerance but not on the specific matrix.

Paper~1 discovered this law empirically on GPT-2 Small (117M
parameters, 72 matrices, $R^2 = 0.90$) and validated it across
GPT-2 Medium (355M, 144 matrices) and Phi-2 (2.7B, 192
matrices)~\citep{calderon2026cartography}.  Paper~2 derived it
analytically: trained weight spectra belong to the stretched
exponential class $\lambda_i \propto e^{-\beta i^\gamma}$, for
which both $\Svn$ and $\log \Dmin$ share the asymptotic scale
$(1/\gamma)|\log \beta|$, forcing $\alpha \to 1$.  Adam
optimization, not SGD, produces the spectral shape that makes the
law hold~\citep{calderon2026theory}.

\subsection{What this paper adds}
\label{sec:intro-contributions}

The first two papers left three questions open.  This paper answers
them.

\begin{enumerate}[leftmargin=2em, itemsep=3pt]
  \item \textbf{Does the law scale?}
    Papers~1 and~2 validated on models up to 2.7B.  Most practical
    compression demand targets the 7B--70B range.  We validate the
    law on Mistral-7B~\citep{jiang2023mistral} (7.2B parameters, 224
    matrices, $R^2 = 0.997$) and Qwen2-7B~\citep{yang2024qwen2}
    (7.1B parameters, 196 matrices, $R^2 = 0.993$).  The fit
    \emph{improves} with scale (\cref{sec:law-scales}).

  \item \textbf{Is the law practically useful?}
    Knowing that $\Svn$ \emph{predicts} compressibility is
    interesting; using that prediction to \emph{allocate} compression
    budgets is useful.  We show that entropy-guided rank allocation
    outperforms uniform allocation at every budget tested on GPT-2
    Small, with margins of 5--38\% in perplexity.  After 200 steps
    of LoRA healing, the advantage grows from 5\% to 11\% at 80\%
    parameter budget (\cref{sec:compression}).

  \item \textbf{How does $\Svn$ compare to existing metrics?}
    AlphaPruning~\citep{lu2024alphapruning} uses the Hill estimator
    $\alpha_{\text{Hill}}$ of spectral tail heaviness to guide
    compression.  On Mistral-7B, $\Svn$ achieves $R^2 = 0.997$ as a
    predictor of $\Dmin$ where $\alpha_{\text{Hill}}$ achieves $R^2 =
    0.022$.  The two metrics are essentially uncorrelated ($r =
    -0.16$); they measure fundamentally different properties of the
    spectrum (\cref{sec:alpha-pruning}).
\end{enumerate}

All experiments run on commodity hardware: spectral analysis of a 7B
model takes approximately 45 minutes on an Apple~M1 with 16\,GB
unified memory and no GPU.  The tool is released as
\texttt{entropy-lens}, an open-source Python package
(\cref{sec:tool}).

\subsection{Outline}
\label{sec:intro-outline}

\Cref{sec:background} reviews SVD compression, the
Entropy-Compression Law, and competing approaches.
\Cref{sec:law-scales} validates the law on Mistral-7B and Qwen2-7B.
\Cref{sec:alpha-pruning} compares $\Svn$ against AlphaPruning's
metric.
\Cref{sec:compression} demonstrates entropy-guided compression with
and without healing.
\Cref{sec:tool} describes the \texttt{entropy-lens} tool.
\Cref{sec:limitations} states our limitations.
\Cref{sec:conclusion} concludes.


% ===========================================================================
% 2. BACKGROUND
% ===========================================================================
\section{Background}
\label{sec:background}

\subsection{SVD compression of weight matrices}
\label{sec:bg-svd}

Given a weight matrix $W \in \mathbb{R}^{m \times n}$, the singular
value decomposition (SVD) writes $W = U \Sigma V^\top$, where
$\Sigma = \operatorname{diag}(\sigma_1, \sigma_2, \ldots)$ with
$\sigma_1 \geq \sigma_2 \geq \cdots \geq 0$.  Truncating at rank
$r$ --- keeping only the $r$ largest singular values --- yields the
best rank-$r$ approximation in the Frobenius
norm~\citep{Eckart1936}.  The \emph{relative Frobenius error} of
this approximation is
\begin{equation}
  \label{eq:frob-error}
  \varepsilon(r) = \sqrt{\frac{\sum_{i=r+1}^{R} \sigma_i^2}
                              {\sum_{i=1}^{R} \sigma_i^2}}\,,
\end{equation}
where $R = \min(m, n)$ is the full rank.  We define $\Dmin(\varepsilon)$
as the smallest $r$ achieving $\varepsilon(r) \leq \varepsilon$ ---
the minimum rank needed for a given error tolerance.

For tensor network compression, the same logic applies at each
internal bond of an MPS decomposition
(\cref{eq:law})~\citep{Oseledets2011, Schollwock2011}.  We use
``bond dimension'' and ``rank'' interchangeably throughout this
paper; both refer to the truncation dimension at a given bipartition.

\subsection{The Entropy-Compression Law}
\label{sec:bg-law}

The full derivation appears in~\citet{calderon2026cartography}
(empirical) and~\citet{calderon2026theory} (analytical).  We
summarize the key ideas.

Reshape a weight matrix $W \in \mathbb{R}^{m \times n}$ into a
high-order tensor by factorizing each dimension into small primes
(\eg, $4096 = 2^{12}$), then decompose the resulting tensor as
an MPS via sequential SVD.  At each internal bipartition --- the
boundary between ``left'' and ``right'' sites --- the SVD produces
a singular value spectrum.  From the normalized squared singular
values $p_i = \sigma_i^2 / \sum_j \sigma_j^2$, we compute the
\emph{von~Neumann entropy} (also called the Shannon entropy of the
squared singular value distribution):
\begin{equation}
  \label{eq:svn}
  \Svn = -\sum_i p_i \log p_i\,.
\end{equation}
This is the standard measure of bipartite entanglement in quantum
information theory~\citep{VonNeumann1932, Eisert2010}.  A flat
spectrum gives $\Svn = \log R$ (maximum entropy, hardest to
compress); a sharply peaked spectrum gives $\Svn \approx 0$ (minimum
entropy, easiest to compress).

The \emph{dominant bipartition} --- the cut with the largest bond
dimension, typically at the boundary between the row and column index
groups of the tensorized matrix --- is the bottleneck for
compression.  The Entropy-Compression Law states that $\Dmin$ at this
cut obeys \cref{eq:law}.  Paper~2 shows that the exponent
$\alpha \to 1$ for the stretched exponential spectral class produced
by Adam training, so that in the asymptotic regime,
$\Dmin \approx c \cdot e^{\Svn}$ --- the bond dimension equals the
exponential of the entropy, which is the effective dimension of the
entangled subspace~\citep{calderon2026theory}.

\subsection{Competing approaches}
\label{sec:bg-competing}

Several recent methods address layer-wise rank allocation for LLM
compression.

\paragraph{AlphaPruning.}
\citet{lu2024alphapruning} use Heavy-Tailed Self-Regularization
(HT-SR) theory to assign compression ratios.  For each weight
matrix, they compute the \emph{Hill estimator}
$\alpha_{\text{Hill}}$, which measures the heaviness of the spectral
tail: a smaller $\alpha_{\text{Hill}}$ indicates a heavier tail (more
weight in large singular values), and AlphaPruning assigns less
compression to matrices with smaller $\alpha_{\text{Hill}}$.  The
metric captures a real property of the spectrum --- its tail behavior
--- but, as we show in \cref{sec:alpha-pruning}, tail heaviness is a
poor predictor of the \emph{total truncation error}, which depends on
the full spectral distribution, not just its tail.

\paragraph{SVD-LLM.}
\citet{Wang2024svd} incorporate truncation-awareness into SVD
compression by whitening the weight matrices before decomposition,
accounting for activation statistics.  Their method improves
compression quality but does not provide a structural predictor of
per-layer compressibility.

\paragraph{Adaptive rank allocation.}
\citet{Sun2025ara} (ARA) propose search-based adaptive rank
allocation, assigning different ranks to different layers based on
perplexity sensitivity.  The Entropy-Compression Law provides the
structural predictor that makes such search unnecessary: $\Svn$ tells
you which layers need more rank, without any evaluation passes.


% ===========================================================================
% 3. THE LAW SCALES
% ===========================================================================
\section{The Law Scales}
\label{sec:law-scales}

\subsection{Experimental setup}
\label{sec:scale-setup}

We apply the Entropy-Compression Law analysis to three models
spanning two orders of magnitude in parameter count:

\begin{itemize}[leftmargin=2em, itemsep=2pt]
  \item \textbf{GPT-2 Small}~\citep{Radford2019}: 124M parameters,
    12 layers, 72 weight matrices, $d_{\text{model}} = 768$.
  \item \textbf{Mistral-7B}~\citep{jiang2023mistral}: 7.2B
    parameters, 32 layers, 224 weight matrices (7 per layer:
    q/k/v/o projections + gate/up/down FFN),
    $d_{\text{model}} = 4096$.
  \item \textbf{Qwen2-7B}~\citep{yang2024qwen2}: 7.1B parameters,
    28 layers, 196 weight matrices (7 per layer),
    $d_{\text{model}} = 3584$.
\end{itemize}

For each model, we extract every weight matrix from the attention and
feed-forward blocks, reshape it into a high-order tensor by
factorizing each dimension into powers of 2 (and one factor of 7 for
Qwen2's 3584-dimensional matrices), decompose it as an MPS via
sequential SVD, and record the full singular value spectrum at the
dominant bipartition.  We then compute $\Svn$ and $\Dmin(\varepsilon)$
for $\varepsilon \in \{5\%, 10\%, 20\%, 50\%\}$.

All spectral analysis runs on a single Apple~M1 processor with
16\,GB unified memory.  No GPU is required.  Processing time is
approximately 45 minutes per 7B model.

\subsection{Validation results}
\label{sec:scale-results}

\Cref{tab:cross-model} reports the Entropy-Compression Law fit
parameters for all three models.

\begin{table}[htbp]
\centering
\caption{Entropy-Compression Law fit:
  $\log \Dmin(\varepsilon) = \alpha \cdot \Svn + \log c$.
  The law holds across all three models and \emph{improves with
  model size}: Mistral-7B achieves $R^2 = 0.997$ at
  $\varepsilon = 50\%$.  The exponent $\alpha$ approaches unity
  at coarse compression, consistent with the theoretical prediction
  of~\citet{calderon2026theory}.}
\label{tab:cross-model}
\begin{tabular}{@{}llccccc@{}}
\toprule
$\varepsilon$ & Model & Params & $n$ & $R^2$ & $\alpha$ \\
\midrule
\multirow{3}{*}{5\%}
  & GPT-2 Small  & 124M & 72  & 0.804 & 0.367 \\
  & Mistral-7B   & 7.2B & 224 & 0.839 & 0.756 \\
  & Qwen2-7B     & 7.1B & 196 & 0.920 & 0.967 \\
\midrule
\multirow{3}{*}{10\%}
  & GPT-2 Small  & 124M & 72  & 0.808 & 0.546 \\
  & Mistral-7B   & 7.2B & 224 & 0.916 & 0.824 \\
  & Qwen2-7B     & 7.1B & 196 & 0.949 & 0.979 \\
\midrule
\multirow{3}{*}{20\%}
  & GPT-2 Small  & 124M & 72  & 0.827 & 0.732 \\
  & Mistral-7B   & 7.2B & 224 & 0.973 & 0.910 \\
  & Qwen2-7B     & 7.1B & 196 & 0.979 & 0.999 \\
\midrule
\multirow{3}{*}{50\%}
  & GPT-2 Small  & 124M & 72  & 0.865 & 1.001 \\
  & Mistral-7B   & 7.2B & 224 & \textbf{0.997} & 1.041 \\
  & Qwen2-7B     & 7.1B & 196 & \textbf{0.993} & 1.020 \\
\bottomrule
\end{tabular}
\end{table}

\begin{figure}[htbp]
\centering
\fbox{\parbox{0.9\columnwidth}{\centering\vspace{3cm}
  \texttt{[Figure: scatter\_3model\_50pct]}\\[4pt]
  Three-model scatter plot of $\log \Dmin$ vs.\ $\Svn$
\vspace{3cm}}}
\caption{\textbf{The Entropy-Compression Law at 7B scale.}
  Log-scale minimum bond dimension $\Dmin$ versus von~Neumann entropy
  $\Svn$ at $\varepsilon = 50\%$ for GPT-2 Small (72 points, blue),
  Mistral-7B (224 points, orange), and Qwen2-7B (196 points, green).
  The 7B models exhibit wider entropy ranges and tighter fits than
  GPT-2 Small.  Mistral-7B achieves $R^2 = 0.997$.}
\label{fig:scatter-3model}
\end{figure}

Three patterns emerge from \cref{tab:cross-model}.

\paragraph{The law improves with scale.}
At $\varepsilon = 50\%$, GPT-2 Small achieves $R^2 = 0.865$;
Mistral-7B reaches $R^2 = 0.997$; Qwen2-7B reaches $R^2 = 0.993$.
This improvement is not accidental.  Larger models have wider
entropy distributions --- the difference between the most
compressible and least compressible layers is larger --- which gives
the regression more variance to explain.  GPT-2 Small's 768-dimensional
matrices have a narrower entropy range than Mistral's 4096-dimensional
matrices, creating a compression ceiling that limits the $R^2$ for
small models.  As model size grows, the law gets \emph{better}, not
worse.

\paragraph{The exponent $\alpha$ converges to unity.}
Across all three models, $\alpha$ increases toward 1 as
$\varepsilon$ increases: from $\alpha \approx 0.37$--$0.97$ at 5\%
error to $\alpha \approx 1.00$--$1.04$ at 50\% error.  At
$\alpha = 1$, the law reduces to $\Dmin \propto e^{\Svn}$ --- the
\emph{effective Hilbert space dimension} --- which is the
information-theoretic prediction for the number of significant
degrees of freedom.  The convergence toward unity from multiple
architectures with different training recipes confirms the
theoretical prediction of~\citet{calderon2026theory}: the
stretched-exponential spectral class forces $\alpha \to 1$ in the
coarse-compression regime.

\paragraph{The law is architecture-agnostic.}
GPT-2 uses multi-head attention with no gating in the FFN.
Mistral-7B uses grouped-query attention (GQA), rotary position
embeddings (RoPE), and a gated FFN with SiLU activation.
Qwen2-7B uses GQA with different group sizes and its own
training recipe.  Despite these architectural differences, the same
exponential relationship holds.  The law depends on the
\emph{spectral structure of trained weight matrices}, not on the
specifics of how those matrices are used during inference.

\subsection{Entanglement heterogeneity at scale}
\label{sec:scale-heterogeneity}

The entanglement structure of Mistral-7B confirms and amplifies the
heterogeneity pattern discovered in GPT-2
Small~\citep{calderon2026cartography}.  \Cref{tab:mistral-heterogeneity}
shows the mean $\Svn$ and $\Dmin$ at $\varepsilon = 50\%$ by
projection type.

\begin{table}[htbp]
\centering
\caption{Entanglement heterogeneity in Mistral-7B.  Attention
  projections (top four rows) have systematically lower entropy and
  smaller bond dimensions than feed-forward layers (bottom three
  rows).  The ratio between the most compressible type
  (\texttt{k\_proj}, $D_{50} = 318$) and the least compressible
  (\texttt{up\_proj}, $D_{50} = 1998$) exceeds $6\times$.}
\label{tab:mistral-heterogeneity}
\begin{tabular}{@{}lccc@{}}
\toprule
Projection type & $\bar{\Svn}$ (nats) & Mean $D_{50}$ & $n$ \\
\midrule
\texttt{k\_proj}    & 6.249 &  318 & 32 \\
\texttt{v\_proj}    & 6.700 &  490 & 32 \\
\texttt{q\_proj}    & 7.156 &  828 & 32 \\
\texttt{o\_proj}    & 7.559 & 1079 & 32 \\
\midrule
\texttt{gate\_proj} & 8.001 & 1875 & 32 \\
\texttt{up\_proj}   & 8.093 & 1998 & 32 \\
\texttt{down\_proj} & 8.080 & 1973 & 32 \\
\bottomrule
\end{tabular}
\end{table}

\begin{figure}[htbp]
\centering
\fbox{\parbox{0.9\columnwidth}{\centering\vspace{3cm}
  \texttt{[Figure: entropy\_map\_mistral7b]}\\[4pt]
  Heatmap of $\Svn$ by layer and projection type
\vspace{3cm}}}
\caption{Entanglement map of Mistral-7B.  Each cell shows $\Svn$
  (in nats) for one weight matrix, organized by layer (horizontal)
  and projection type (vertical).  The systematic separation between
  attention types (cooler colors) and FFN types (warmer colors) is
  visible at a glance.  Key projections are consistently the most
  compressible; gate and up projections the least.}
\label{fig:entropy-map}
\end{figure}

The separation between attention and FFN layers is decisive:
$\bar{\Svn}^{\text{attn}} = 6.916\nats$ versus
$\bar{\Svn}^{\text{ffn}} = 8.058\nats$ (two-sample $t$-test,
$p = 3.26 \times 10^{-40}$).  This means attention projections
require, on average, roughly $e^{8.06 - 6.92} = e^{1.14} \approx
3.1\times$ less bond dimension than FFN projections at
$\varepsilon = 50\%$.

Within attention, the ordering is k < v < q < o, which differs from
GPT-2's ordering (o < q $\approx$ k < v).  This difference likely
reflects architectural choices: Mistral uses grouped-query attention
where key and value projections have smaller output dimensions
($d_{\text{kv}} = 1024$ vs.\ $d_{\text{model}} = 4096$), making
them inherently lower-rank and lower-entropy.  The Entropy-Compression
Law captures this difference automatically --- no architectural
knowledge is needed, just the SVD of the weight matrix.


% ===========================================================================
% 4. S1 VS ALPHAPRUNING
% ===========================================================================
\section{$\Svn$ vs.\ AlphaPruning}
\label{sec:alpha-pruning}

\subsection{Head-to-head comparison}
\label{sec:alpha-head}

AlphaPruning~\citep{lu2024alphapruning} uses the Hill estimator
$\alpha_{\text{Hill}}$ --- a measure of how heavy-tailed the singular
value distribution is --- to guide layer-wise compression.  The
intuition is that matrices with heavier tails (smaller
$\alpha_{\text{Hill}}$) concentrate more energy in their top singular
values and should therefore be easier to compress.

We computed both $\Svn$ and $\alpha_{\text{Hill}}$ for all 224 weight
matrices of Mistral-7B and used each as a predictor of
$\log \Dmin(\varepsilon)$ via ordinary least squares.
\Cref{tab:s1-vs-alpha} reports the results.

\begin{table}[htbp]
\centering
\caption{$\Svn$ vs.\ $\alpha_{\text{Hill}}$ as predictors of
  $\log \Dmin$ on Mistral-7B ($n = 224$).  At every error tolerance,
  $\Svn$ explains more than an order of magnitude more variance.
  At $\varepsilon = 50\%$, $\alpha_{\text{Hill}}$ explains only 2.2\%
  of the variance in bond dimension --- essentially noise.}
\label{tab:s1-vs-alpha}
\begin{tabular}{@{}lcccc@{}}
\toprule
$\varepsilon$ & $R^2(\Svn)$ & $R^2(\alpha_{\text{Hill}})$ &
  $\Delta R^2$ \\
\midrule
 5\% & 0.839 & 0.100 & $+0.739$ \\
10\% & 0.916 & 0.078 & $+0.838$ \\
20\% & 0.973 & 0.054 & $+0.920$ \\
50\% & 0.997 & 0.022 & $+0.975$ \\
\bottomrule
\end{tabular}
\end{table}

\begin{figure}[htbp]
\centering
\fbox{\parbox{0.9\columnwidth}{\centering\vspace{3cm}
  \texttt{[Figure: s1\_vs\_alpha\_hill]}\\[4pt]
  Side-by-side scatter: $\Svn$ vs.\ $\alpha_{\text{Hill}}$ as predictors
\vspace{3cm}}}
\caption{Scatter plots of $\Svn$ (left) and $\alpha_{\text{Hill}}$
  (right) versus $\log \Dmin$ at $\varepsilon = 50\%$ for all 224
  Mistral-7B matrices.  $\Svn$ produces a tight linear relationship
  ($R^2 = 0.997$); $\alpha_{\text{Hill}}$ shows no meaningful
  pattern ($R^2 = 0.022$).}
\label{fig:s1-vs-alpha}
\end{figure}

The comparison is not close.  At $\varepsilon = 50\%$, $\Svn$
explains 99.7\% of the variance in $\log \Dmin$;
$\alpha_{\text{Hill}}$ explains 2.2\%.  At $\varepsilon = 5\%$, where
the fit is weakest, $\Svn$ still explains 83.9\% versus
$\alpha_{\text{Hill}}$'s 10.0\%.

\subsection{Why the metrics diverge}
\label{sec:alpha-why}

The Pearson correlation between $\Svn$ and $\alpha_{\text{Hill}}$
across Mistral-7B's 224 matrices is $r = -0.16$ --- the two metrics
are essentially uncorrelated.  They measure different things.

$\Svn$ measures the \emph{full spectral distribution}.  It is the
Shannon entropy of the squared singular value probabilities
$\{p_i\}$, giving weight to every part of the spectrum through the
$-p_i \log p_i$ term.  The truncation error
(\cref{eq:frob-error}) is the sum of all discarded $\sigma_i^2$
values --- head, body, and tail combined --- so $\Svn$, which
accounts for the entire distribution, is the natural predictor.

$\alpha_{\text{Hill}}$ measures the \emph{tail heaviness} of the
spectrum.  It is estimated from the $k$ largest singular values and
characterizes how quickly the tail decays.  A matrix can have a heavy
tail (small $\alpha_{\text{Hill}}$) but still require a large bond
dimension because its \emph{body} --- the middle portion of the
spectrum --- carries significant energy.  Conversely, a matrix with a
light tail (large $\alpha_{\text{Hill}}$) might need few modes
overall because its energy is concentrated in the top singular values.

The $r = -0.16$ correlation confirms this: knowing the tail heaviness
tells you almost nothing about the total truncation error.  It is as
if you tried to predict a river's total flow by measuring only its
depth at the deepest point --- you would miss the width entirely.
$\Svn$ measures the whole cross-section.


% ===========================================================================
% 5. ENTROPY-GUIDED COMPRESSION
% ===========================================================================
\section{Entropy-Guided Compression}
\label{sec:compression}

\subsection{The allocation problem}
\label{sec:comp-problem}

Given a global \emph{parameter budget} --- the total number of
parameters allowed after compression, expressed as a fraction of the
original model --- we must decide how many parameters to allocate to
each layer.  This is an optimization problem: minimize total
approximation error (or perplexity degradation) subject to a global
parameter constraint.

A weight matrix $W \in \mathbb{R}^{m \times n}$ truncated at rank
$r$ uses $r(m + n)$ parameters (storing the factors $U_r$ and
$V_r^\top$).  The parameter \emph{savings} from compressing one
matrix are $(mn - r(m+n))$, and the total savings across all
matrices must reach the target budget.

\subsection{Three strategies}
\label{sec:comp-strategies}

We compare three rank allocation strategies, all subject to the same
global parameter budget:

\begin{enumerate}[leftmargin=2em, itemsep=3pt]
  \item \textbf{Uniform.}
    Every matrix receives the same truncation rank, computed so that
    the total parameter count matches the budget.  This is the
    baseline strategy used by most SVD compression methods.

  \item \textbf{Proportional.}
    Each matrix's rank is proportional to its original size.  Larger
    matrices get proportionally more rank.  This preserves the
    relative complexity of each layer but ignores their differing
    tolerances for compression.

  \item \textbf{Entropy-guided.}
    We use the Entropy-Compression Law to predict
    $\Dmin(\varepsilon)$ for each matrix at a reference $\varepsilon$,
    then rescale all ranks proportionally to fit the global budget.
    Matrices with high $\Svn$ (harder to compress) receive more rank;
    matrices with low $\Svn$ (easier to compress) receive less.  The
    $\Svn$ values are computed once, offline, and the allocation
    requires no evaluation passes.
\end{enumerate}

\subsection{Results on GPT-2 Small}
\label{sec:comp-results}

We compressed GPT-2 Small at six parameter budgets (60\%--95\% of
original parameters retained) using uniform and entropy-guided
allocation, then evaluated perplexity on WikiText-2.
\Cref{tab:compression} reports the results.

\begin{table}[htbp]
\centering
\caption{Perplexity (WikiText-2) of GPT-2 Small under SVD compression
  with uniform vs.\ entropy-guided rank allocation.  Baseline
  (uncompressed) PPL is 31.83.  Entropy-guided allocation is better
  at \emph{every} budget tested, with advantages ranging from 5\% to
  38\%.  Lower perplexity is better.}
\label{tab:compression}
\begin{tabular}{@{}rrrr@{}}
\toprule
Budget & Uniform PPL & Entropy PPL & Improvement \\
\midrule
95\% & 4783 & 3016 & 37\% \\
90\% & 5505 & 3401 & 38\% \\
85\% & 4397 & 3491 & 21\% \\
80\% & 3965 & 3751 & ~5\% \\
70\% & 4468 & 4263 & ~5\% \\
60\% & 7360 & 6270 & 15\% \\
\bottomrule
\end{tabular}
\end{table}

\begin{figure}[htbp]
\centering
\fbox{\parbox{0.9\columnwidth}{\centering\vspace{3cm}
  \texttt{[Figure: pareto\_compression]}\\[4pt]
  PPL vs.\ budget for uniform and entropy-guided strategies
\vspace{3cm}}}
\caption{Pareto frontier: perplexity versus parameter budget for
  uniform and entropy-guided compression of GPT-2 Small.
  Entropy-guided allocation dominates the Pareto frontier at every
  budget.  The advantage is largest at aggressive compression (90\%
  budget: 38\% improvement) and smallest at moderate compression
  (80\% budget: 5\% improvement).}
\label{fig:pareto}
\end{figure}

Entropy-guided allocation dominates at every budget.  The advantage
is largest at the most aggressive compression ratios (90\% and 95\%
budget, where 5--10\% of parameters are removed).  At these ratios,
the uniform strategy is forced to truncate every matrix by the same
small amount, including high-entropy FFN matrices that cannot
tolerate it.  The entropy-guided strategy instead concentrates the
cuts on low-entropy attention matrices that can absorb them with
minimal damage.

We note that the absolute perplexity values are high (thousands,
versus a baseline of 31.83) because raw SVD truncation without any
post-compression fine-tuning is an aggressive intervention.  The
relevant comparison is not absolute perplexity but the
\emph{relative} performance of the two allocation strategies.  The
next section shows that healing closes most of the gap to baseline.

\subsection{Healing amplifies the advantage}
\label{sec:comp-healing}

Raw SVD truncation degrades perplexity severely because the model
has never seen its own compressed weights during training.
\emph{Healing} --- a short fine-tuning pass on the compressed model
--- allows the surviving parameters to adapt to the new weight
configuration and recover much of the lost performance.

We heal both the uniform and entropy-guided compressed models at
80\% parameter budget using LoRA~\citep{Hu2021} (rank 16,
$\alpha = 32$, learning rate $2 \times 10^{-4}$) for 200 steps on
WikiText-2.  \Cref{tab:healing} reports the results.

\begin{table}[htbp]
\centering
\caption{Healing amplifies the entropy advantage.  At 80\% parameter
  budget, entropy-guided compression starts 5\% better before healing
  and 11\% better after healing.  The advantage nearly doubles because
  healing is more effective when the initial compression preserves the
  most critical parameters.  Baseline PPL: 31.83.}
\label{tab:healing}
\begin{tabular}{@{}lccc@{}}
\toprule
Strategy & Compressed PPL & Healed PPL & Improvement \\
\midrule
Uniform        & 3965  & 209.89 & --- \\
Entropy-guided & 3751  & 186.40 & 11\% \\
\bottomrule
\end{tabular}
\end{table}

After healing, the entropy-guided model achieves PPL~186.40 versus
the uniform model's 209.89 --- an 11\% improvement, up from the
5\% advantage before healing.  The advantage nearly doubles because
healing is more effective when the initial compression preserves the
right parameters.  Entropy-guided allocation concentrates rank where
the spectrum demands it, giving the LoRA adapter a better starting
point for recovery.

This result has a practical implication: in any compression pipeline
that includes a healing step (which all production pipelines should),
the benefit of entropy-guided allocation is \emph{amplified}, not
diminished.  Spending 45~minutes to compute $\Svn$ before
compressing saves far more than 45~minutes of additional healing
compute.


% ===========================================================================
% 6. ENTROPY-LENS: A PRACTICAL TOOL
% ===========================================================================
\section{\texttt{entropy-lens}: A Practical Tool}
\label{sec:tool}

The value of the Entropy-Compression Law depends on how cheaply
$\Svn$ can be computed.  If measuring the entropy of a 7B model
required a GPU cluster, the law would be an interesting theoretical
result but not a practical one.  We show that it is practical.

\subsection{Architecture}
\label{sec:tool-arch}

\texttt{entropy-lens} is a Python package that computes $\Svn$ for
every weight matrix of any HuggingFace model.  The key design
decisions are:

\begin{enumerate}[leftmargin=2em, itemsep=3pt]
  \item \textbf{Streaming SVD via safetensors mmap.}
    Instead of loading the entire model into memory, we memory-map
    the safetensors files and process one weight matrix at a time.
    Each matrix is loaded, tensorized, SVD-decomposed at the dominant
    bipartition, and its spectrum recorded.  Peak memory usage is
    dominated by the single largest weight matrix (typically
    $d_{\text{model}} \times 4 d_{\text{model}}$ for the FFN),
    not by the total model size.

  \item \textbf{Architecture-agnostic extraction.}
    A registry maps HuggingFace model families (GPT-2, LLaMA/Mistral,
    Qwen, Phi) to their weight matrix naming conventions.  Adding a
    new architecture requires implementing a single
    \texttt{ArchConfig} subclass that specifies the projection names
    and their shapes.

  \item \textbf{Dominant-bipartition focus.}
    Rather than computing the full MPS decomposition (which would
    require $N-1$ SVDs per matrix), we compute a single SVD at the
    dominant bipartition --- the cut with the largest dimension,
    which is the compression bottleneck.  This reduces the per-matrix
    cost from $O(N)$ SVDs to one.
\end{enumerate}

\subsection{Performance}
\label{sec:tool-perf}

\Cref{tab:tool-perf} reports the wall-clock time for analyzing
three models on an Apple~M1 with 16\,GB unified memory.

\begin{table}[htbp]
\centering
\caption{\texttt{entropy-lens} performance on Apple~M1 (16\,GB, no
  GPU).  Analyzing a 7B model takes under 45 minutes.  The dominant
  cost is the SVD of the largest weight matrices ($4096 \times 14336$
  for Mistral's FFN layers).}
\label{tab:tool-perf}
\begin{tabular}{@{}lccc@{}}
\toprule
Model & Params & Matrices & Time \\
\midrule
GPT-2 Small & 124M &  72 & $\sim$2 min \\
Mistral-7B  & 7.2B & 224 & $\sim$45 min \\
Qwen2-7B    & 7.1B & 196 & $\sim$40 min \\
\bottomrule
\end{tabular}
\end{table}

\subsection{Open-source release}
\label{sec:tool-release}

\texttt{entropy-lens} is released under the MIT license at
\url{https://github.com/nicogarazzo/entropy-lens}.  The tool provides
a command-line interface:

\begin{verbatim}
  entropy-lens analyze mistralai/Mistral-7B-v0.3
\end{verbatim}

\noindent
This produces a CSV of per-matrix $\Svn$, $\Dmin(\varepsilon)$ at
four thresholds, scatter plots of the Entropy-Compression Law fit,
and an entanglement heatmap.  The results can be fed directly into
the rank allocation pipeline.


% ===========================================================================
% 7. LIMITATIONS
% ===========================================================================
\section{Limitations}
\label{sec:limitations}

We state our limitations directly.

\begin{enumerate}[leftmargin=2em, itemsep=3pt]
  \item \textbf{Validated up to 7B, not 70B+.}
    The largest model tested is 7.2B parameters.  Whether the law
    holds at the 70B--405B scale where compression is most urgently
    needed remains to be validated.  We expect it will --- the law
    \emph{improved} from 124M to 7.2B --- but expectation is not
    evidence.  \texttt{entropy-lens} can analyze 70B models in
    principle (via streaming SVD), but we have not yet run this
    experiment.

  \item \textbf{Frobenius error $\neq$ perplexity.}
    The Entropy-Compression Law predicts Frobenius truncation error,
    not perplexity.  As~\citet{hsu2025rethinking} emphasize, the
    relationship between per-layer approximation error and downstream
    task performance is non-trivial.  Our compression experiments
    (\cref{sec:compression}) show that the Frobenius-based ranking
    of layers is informative for perplexity-aware allocation, but we
    have not derived a formal perplexity-prediction law.

  \item \textbf{Healing is obligatory for practical compression.}
    Without post-compression fine-tuning, SVD truncation produces
    perplexity in the thousands, regardless of allocation strategy.
    The entropy-guided advantage is in \emph{relative} performance
    (which strategy degrades less), not in absolute usability of the
    compressed model without healing.  Any practical deployment
    requires a healing step.

  \item \textbf{Compression experiments on GPT-2 Small only.}
    The perplexity comparison between uniform and entropy-guided
    allocation (\cref{sec:compression}) is currently validated only
    on GPT-2 Small.  Mistral-7B compression experiments require GPU
    cloud infrastructure and are in progress.  We include the GPT-2
    results because the allocation principle --- entropy predicts
    which layers need more rank --- is architecture-independent, even
    though the specific perplexity numbers are model-specific.

  \item \textbf{Single bipartition.}
    We measure $\Svn$ at the dominant bipartition only.  A full
    entanglement characterization would examine all $N-1$ bipartitions
    and their mutual information structure.  For the purpose of
    compression, the dominant bipartition is the bottleneck, so this
    is the relevant cut --- but the full multi-cut structure may
    contain additional information for more sophisticated compression
    schemes.

  \item \textbf{No comparison against perplexity-aware search methods.}
    We compare entropy-guided allocation against uniform allocation,
    not against search-based methods like ARA~\citep{Sun2025ara} that
    optimize ranks by evaluating perplexity at many configurations.
    Such methods may find better allocations at the cost of orders of
    magnitude more compute.  The entropy-guided approach is intended
    as a \emph{principled initialization} --- a starting point that
    eliminates most of the search space, not necessarily the final
    answer.
\end{enumerate}


% ===========================================================================
% 8. CONCLUSION
% ===========================================================================
\section{Conclusion}
\label{sec:conclusion}

The Entropy-Compression Law scales.  On Mistral-7B, the von~Neumann
entanglement entropy $\Svn$ predicts the minimum bond dimension for
a given Frobenius error with $R^2 = 0.997$ --- explaining 99.7\% of
the variance across 224 weight matrices with a single structural
measurement per matrix.  On Qwen2-7B, the law achieves $R^2 = 0.993$.
Both results are stronger than the original GPT-2 validation
($R^2 = 0.865$), because larger models exhibit wider entropy
variation across layers, giving the law more leverage.

Beyond prediction, the law is practically useful.  Entropy-guided
rank allocation outperforms uniform allocation at every compression
budget tested, with advantages of 5--38\% in perplexity before
healing and 11\% after healing.  As a predictor of layer-wise
compressibility, $\Svn$ dominates AlphaPruning's Hill estimator
($R^2 = 0.997$ vs.\ $0.022$); the two metrics are essentially
independent ($r = -0.16$), measuring different properties of the
singular value spectrum.  All computations run on a laptop in under
an hour.

One open question is whether the gap between Frobenius prediction and
perplexity can be closed --- whether a perplexity-aware version of
the Entropy-Compression Law exists.  Another is whether the law
holds at 70B+ scale.  But the direction of the evidence is clear:
the law gets better with scale, not worse.  If the trend continues,
the largest models --- the ones that need compression most --- are
the ones where entropy-guided allocation will work best.

A quantum-inspired diagnostic, requiring no labels, no gradients,
and no GPU, predicts the compressibility of every layer in a 7B
language model with 99.7\% accuracy.  The singular value spectrum
of a trained weight matrix, it turns out, already knows how much
you can compress it.  You just have to ask it the right question.


% ===========================================================================
% ACKNOWLEDGMENTS
% ===========================================================================
\section*{Acknowledgments}

Computations were performed on Apple~Silicon~M1 hardware.  The
author thanks the open-source communities behind
PyTorch~\citep{Paszke2019} and HuggingFace
Transformers~\citep{Wolf2020}.

\paragraph{AI-assisted research statement.}
This research was conducted by a human author with substantial
assistance from large language models (Anthropic's Claude), used for
code development, literature review, and manuscript drafting.  All
experiments were designed, executed, and verified by the author; all
numerical results come from the \texttt{entropy-lens} code
repository, which can be run end-to-end to reproduce them.  The
author takes full responsibility for every claim.


% ===========================================================================
% BIBLIOGRAPHY
% ===========================================================================
\bibliographystyle{plainnat}
\bibliography{references}


\end{document}
